% Direct-authority diplomatic French for printed p.197.
\clearpage
\SourcePage{197}

\LayerTarget{tohoku:conclusion:injective-g-o-embedding:p1}
suite d'obtenir une forme explicite commode pour ``suffisamment'' d'objets
injectifs, en généralisant la construction faite dans 3.1.
\LayerTarget{tohoku:construction:pointwise-g-o-action}
Soit $(A_x)_{x\in X}$ une famille de ${\mathbf O}(x)$-modules $A_x$, considérons le
${\mathbf O}$-Module produit qu'elle définit (cf. 3.1). Pour y faire opérer $G$ de façon
à obtenir un $G$-${\mathbf O}$-faisceau, il suffit pour tout $x\in X$ et $g\in G$ de se
donner un homomorphisme $a\to g.\,a$ du groupe abélien
$A_x$ dans le groupe abélien $A_{g.\, x}$, de façon qu'on ait $g.\,(ua)=(g.\,u)(g.\,a)$ pour $u\in {\mathbf O}(x)$ et $a\in A_x$, et que $e.\,a=a$, $g.\,(g'.\,a)=(gg').\,a$.
\LayerTarget{tohoku:construction:auxiliary-ring-u-x}
Introduisant alors l'anneau auxiliaire $U_x$ engendré par l'algèbre
$\mathbf Z(G_x)$ de $G_x$ et l'anneau ${\mathbf O}(x)$, soumis aux relations de
commutation $gug^{-1}=u^g$ pour $u\in {\mathbf O}(x)$, $g\in G_x$
(où on note $u^g$, pour éviter des confusions, le transformé de $u$ par
$g\in G_x$), on voit que les $A_x$ sont en fait des $U_x$-modules, et un
$g\in G$ quelconque définit un isomorphisme
$a\to g.\,a$ du groupe $A_x$ sur $A_{g.\, x}$
satisfaisant à $g.\,(ua)=(g.\,u)(g.\,a)$ pour $u\in U_x$, $a\in A_x$.
\LayerTarget{tohoku:construction:orbit-representative-u-y}
Il en résulte facilement que, si on choisit un élément $\xi(y)$ dans toute
trajectoire $y\in Y$, et si on pose $U_y=U_{\xi(y)}$, alors les données
précédentes sont équivalentes à la donnée d'une famille $(A_y)_{y\in Y}$ de
$U_y$-modules $A_y$ $(=A_{\xi(y)})$.

\LayerTarget{tohoku:construction:induced-a-bar-y}
Les $A_x$ ($x\in y$) peuvent se déduire de $A_y$ de la façon suivante. Soit
$\overline A_y$ le $G$-module \textit{induit} par le $G_y$-module $A_y$, i.e.
le $G$-module $\mathrm{Hom}_{G_y}(\mathbf Z(G),A_y)$ obtenu à partir du
$G_y$-module $A_y$ par ``extension contravariante des scalaires'' : $A_y$
s'identifie alors à un groupe quotient de $\overline A_y$ par un sous-groupe
abélien stable par $G_y$, et $\overline A_y$ s'identifie au produit de tout
les quotients transformés $g.\,A_y$ pour $g\in G/G_y$.
\LayerTarget{tohoku:construction:induced-o-bar-y}
Introduisant de même le $G$-module $\overline O_y$ induit par
${\mathbf O}_y={\mathbf O}(\xi(y))$, on constate aussitôt que $\overline A_y$ est un
$\overline O_y$-module s'identifiant au produit des
$g.\,O_y$-modules $g.\,A_y$ pour $g\in G/G_y$. On a
alors des isomorphismes canoniques $g.\,{\mathbf O}_y={\mathbf O}(g.\,\xi(y))$, $g.\,A_y=A_{g.\,\xi(y)}$.
\LayerTarget{tohoku:construction:p-a-g-o-sheaf}
Désignons par $P(A)$ le $G$-${\mathbf O}$-faisceau sur $X$ défini par la famille
$A=(A_y)$ de $U_y$-modules $A_y$.
\LayerTarget{tohoku:comparison:hom-to-p-a}
Soit $B$ un $G$-${\mathbf O}$-faisceau quelconque. Un ${\mathbf O}$-homomorphisme d'un
$G$-${\mathbf O}$-Module $B$ dans $P(A)$ s'identifie à une famille $(v_x)_{x\in X}$ de
${\mathbf O}(x)$-homomorphismes $B(x)\to A_x$ (comme nous l'avions signalé
dans 3.1), et ce ${\mathbf O}$-homomorphisme est compatible avec $G$ si et seulement si
les $v_x$ ``permutent'' à $G$ ; on voit alors tout de suite qu'il revient au
même de se donner un $G$-${\mathbf O}$-homomorphisme de $B$ dans $P(A)$, où une famille
$(v_y)_{y\in Y}$ de $U_y$-homomorphismes
$B(\xi(y))\to A_y$. On en conclut aussitôt :

\LayerTarget{tohoku:proposition:5.1.2}
\textsc{Proposition 5.1.2.} \textit{Soit $(A_y)_{y\in Y}$ une
famille de $U_y$-modules injectifs, alors le $G$-${\mathbf O}$-Moduie $P((A_y))$
qu'elle définit est injectif. De plus, tout $G$-${\mathbf O}$-Module est isomorphe à un
sous-faisceau d'un faisceau du type précédent.}

\LayerTarget{tohoku:proof:5.1.2}
Ce dernier fait (le plus important pour nous) est immédiat, en plongeant pour
tout $y$, le module ponctuel $B(\xi(y))$ du faisceau donné $B$ dans un
$U_y$-module injectif.

\LayerTarget{tohoku:construction:direct-image-g-sheaf:p1}
\textit{Images directes de $G$-faisceaux.} Soit $A$ un faisceau
d'ensembles sur $X$, rappelons qu'on a défini son image directe $f_*(A)$ comme
le faisceau sur $Y$ dont les sections sur l'ouvert $U\subset Y$ sont les
sections de $A$ sur $f^{-1}(U)$. Si $A$ est un $G$-faisceau, alors
$\Gamma(f^{-1}(U),A)$ admet $G$ comme groupe d'opérateurs, d'où on conclut
facilement que \textit{l'image directe $f_*(A)$ d'un $G$-faisceau sur $X$ est
un $G$-faisceau sur $Y$} (rappelons que $G$ opère trivialement sur $Y$) ; bien

\TohokuContinuation{198}
